\scanpage{159}
de la forme $d^0_E(u({\unsure{\xi|\zeta}}'))$, $\xi'\in E'(\cdot)$, on voit que $\Delta$ définit une section
\[
 \unsure{\nu|v}_\Delta(P',d_{P'})\in\Gamma(S,F). \tag*{(46)}
\]
Ici la situation universelle est obtenue en prenant pour $F$ le préfaisceau
\[
 F=P^\infty_{S/T}(E\otimes\Omega^{1}_{S/T})/\operatorname{Im}P^\infty_{S/T}(E')\,,
\]
où $P^\infty_{S/T}(E')\longrightarrow P^\infty_{S/T}(E\otimes\Omega_{S/T})$ provient de l'op. diff. composé (de \uline{degré} 1) \inserted{$d^0_Eu$:} $E'\mathrel{\overset{u}{\text{---}}}E\mathrel{\overset{d^0_E}{\text{---}}}E\otimes\Omega_{S/T}$.

Revenons au cas où $E$ admet une stratification, \deleted{et $P$ muni d'une stratification, compatible avec} et \unsure{prenant|prenons} $d^0_E$ déduit d'icelle ; \inserted{soit $(P',\deleted{\unsure{hom}}\sigma)$} \deleted{pour ramener} (comme on \unsure{disait|didait} \unsure{au|que} \unsure{n°|où}. \deleted{les sections de la forme} On ramène les sections $\nu_\Delta$ aux $\mu_D$ de la façon suivante. À tout \deleted{\unsure{$\nu_\Delta$}} \inserted{$\Delta$} est associé un opérateur de Picard-Fuchs, \deleted{\unsure{savoir pour la}} pour (30), en prenant le composé
\[
 D_0:E\xrightarrow{d^0_E}E\otimes\Omega_{S/T}\xrightarrow{\Delta}F \,, \tag*{(47)}
\]
et en passant au quotient, \deleted{\unsure{ce qui}} \inserted{en $D:E''\longrightarrow F$,} ce qui est possible puisque par hyp. sur $\Delta$\quad $\Delta\circ d^0_E$ est nulle sur $E'$. Ceci posé, on trouve que
\[
 \nu_\Delta=\mu_D \,, \tag*{(48)}
\]
% slanted
\margin{faire la vérification}
ce qui donne la réduction annoncée.
% archive sheet number: 158
